\chapter{Manual Override Tools}
\label{chap:Manual_Override_Tools}

\section{Purpose}

Manual Override Tools provide emergency mechanisms for human operators to:
\begin{itemize}
    \item halt dangerous behaviour,
    \item restrict CME capabilities,
    \item roll back recent changes,
    \item isolate or quarantine Engrams.
\end{itemize}

They are designed for rare use when automated and policy-level safeguards
are insufficient.

\section{Override Types}

\subsection{Soft Overrides}

\begin{itemize}
    \item temporarily suspend a specific SIGIL Pathway,
    \item pause canonicalization of Engrams,
    \item dampen access to certain domains or roles.
\end{itemize}

\subsection{Hard Overrides}

\begin{itemize}
    \item quarantine large subsets of Engrams or memory,
    \item revoke keys or access scopes,
    \item stop specific deployment endpoints.
\end{itemize}

\subsection{Emergency Halt}

\begin{itemize}
    \item shut down the CME’s active instances,
    \item freeze Engram writes to C- and P-class,
    \item require EAC review before resumption.
\end{itemize}

\section{Invocation Conditions}

Manual overrides SHOULD only be used when:

\begin{itemize}
    \item a clear safety or misalignment risk is present,
    \item SIGIL or normal CAR processes are too slow or unavailable,
    \item incident policies declare a relevant trigger.
\end{itemize}

Use of Emergency Halt MUST be:
\begin{itemize}
    \item rare,
    \item well-justified,
    \item immediately logged,
    \item followed by formal EAC and CAR review.
\end{itemize}

\section{Logging and Accountability}

Every manual override action MUST:
\begin{itemize}
    \item be recorded in the HITL audit log,
    \item include operator identity and reason,
    \item document collateral impact and follow-up steps.
\end{itemize}

\section{Summary}

Manual Override Tools are the last-resort levers of control. They exist to
protect humans and systems when other lines of defence are insufficient,
and must be wielded with care and accountability.
